% =====================================================================
%  preamble.tex -- packages, styles and the flowchart visual language
%  Smart Agriculture Digital Twin -- KUKA YouBot -- PFA technical report
% =====================================================================
%
%  Everything that is not prose lives here so the section files stay
%  readable. Three things are defined below and used everywhere else:
%
%    1. IEEE-conformant typography (IEEEtran, Times, two columns).
%    2. A code-listing style per language (bash / python / xml / yaml).
%    3. THE FLOWCHART LANGUAGE: a fixed set of coloured node shapes and
%       three phase bands -- INITIALISATION, MAIN LOOP, TERMINATION --
%       so that every flowchart in this report reads the same way.
%
%  Compile with:  make        (see Makefile, uses pdflatex + bibtex)
% =====================================================================

% ------------------------------------------------------- encoding, fonts
\usepackage[T1]{fontenc}
\usepackage[utf8]{inputenc}
\usepackage[english]{babel}
\usepackage{textcomp}

% ------------------------------------------------------- maths
\usepackage{amsmath}
\usepackage{amssymb}
\usepackage{amsfonts}
\usepackage{bm}

% ------------------------------------------------------- tables & figures
\usepackage{graphicx}
\usepackage{booktabs}
\usepackage{array}
\usepackage{multirow}
\usepackage{tabularx}
\usepackage{longtable}
\usepackage{float}
\usepackage{caption}
\usepackage{subcaption}

% ------------------------------------------------------- code listings
\usepackage{listings}
\usepackage{xcolor}

% ------------------------------------------------------- pseudocode
\usepackage{algorithm}
\usepackage{algpseudocode}

% ------------------------------------------------------- drawings
\usepackage{tikz}
\usetikzlibrary{shapes.geometric,shapes.misc,shapes.symbols,arrows.meta,
                positioning,fit,backgrounds,calc,decorations.pathreplacing,
                patterns}

% ------------------------------------------------------- IEEE citations
\usepackage{cite}
\usepackage{url}

% ------------------------------------------------------- misc
\usepackage{enumitem}
\usepackage{ifthen}

% hyperref must come last
\usepackage[hidelinks,breaklinks=true]{hyperref}

% =====================================================================
%  1. COLOUR PALETTE
% =====================================================================
%  One palette for the whole document. The three phase colours are the
%  spine of every flowchart; the accent colours are for emphasis inside
%  tables and listings. Chosen to stay legible when printed in grey.
% ---------------------------------------------------------------------

% --- flowchart phase bands (light fills, saturated borders)
\definecolor{phaseInitFill}{HTML}{DCE9F7}   % blue   -- INITIALISATION
\definecolor{phaseInitLine}{HTML}{2E5F9E}
\definecolor{phaseMainFill}{HTML}{E2F3E4}   % green  -- MAIN LOOP
\definecolor{phaseMainLine}{HTML}{2E7D32}
\definecolor{phaseEndFill}{HTML}{FBE4DC}    % red    -- TERMINATION
\definecolor{phaseEndLine}{HTML}{C0392B}

% --- flowchart node fills
\definecolor{nodeTerm}{HTML}{FFD9A0}        % terminal  (start / stop)
\definecolor{nodeTermLine}{HTML}{B26B00}
\definecolor{nodeProc}{HTML}{FFFFFF}        % process
\definecolor{nodeProcLine}{HTML}{37474F}
\definecolor{nodeDec}{HTML}{FFF6C2}         % decision
\definecolor{nodeDecLine}{HTML}{B8860B}
\definecolor{nodeIO}{HTML}{DDE6FA}          % input / output (ROS topic)
\definecolor{nodeIOLine}{HTML}{3F51B5}
\definecolor{nodeSub}{HTML}{E8DAF5}         % subroutine / library call
\definecolor{nodeSubLine}{HTML}{6A1B9A}
\definecolor{nodeNote}{HTML}{F5F5F5}        % annotation
\definecolor{nodeFail}{HTML}{FADBD8}        % failure / recovery branch
\definecolor{nodeFailLine}{HTML}{922B21}

% --- semantic accents used in prose and tables
\definecolor{okGreen}{HTML}{1E7B34}
\definecolor{warnAmber}{HTML}{B7791F}
\definecolor{badRed}{HTML}{B03A2E}
\definecolor{codeGray}{HTML}{4F4F4F}
\definecolor{codeBg}{HTML}{F7F7F9}
\definecolor{codeKey}{HTML}{0B5394}
\definecolor{codeStr}{HTML}{A31515}
\definecolor{codeCom}{HTML}{5B7C5B}

% shorthand for verdict marks in SWOT and results tables
\newcommand{\ok}[1]{\textcolor{okGreen}{\textbf{#1}}}
\newcommand{\warn}[1]{\textcolor{warnAmber}{\textbf{#1}}}
\newcommand{\bad}[1]{\textcolor{badRed}{\textbf{#1}}}

% =====================================================================
%  2. CODE LISTINGS
% =====================================================================
%  Two columns are narrow, so listings default to \footnotesize and
%  break long lines with a visible arrow. Wide listings go in a
%  figure*/table* float instead of being shrunk further.
% ---------------------------------------------------------------------
\lstdefinestyle{base}{
  basicstyle       = \ttfamily\footnotesize\color{codeGray},
  keywordstyle     = \color{codeKey}\bfseries,
  stringstyle      = \color{codeStr},
  commentstyle     = \color{codeCom}\itshape,
  numberstyle      = \tiny\color{gray},
  backgroundcolor  = \color{codeBg},
  frame            = single,
  rulecolor        = \color{gray!40},
  framesep         = 3pt,
  xleftmargin      = 6pt,
  xrightmargin     = 2pt,
  breaklines       = true,
  breakatwhitespace= false,
  postbreak        = \mbox{\textcolor{gray}{$\hookrightarrow$}\space},
  showstringspaces = false,
  tabsize          = 2,
  columns          = fullflexible,
  keepspaces       = true,
  upquote          = true,
  captionpos       = b,
}
\lstdefinestyle{bash}{style=base,language=bash,
  morekeywords={ros2,colcon,source,export,sudo,apt,git,gz,rviz2,make}}
\lstdefinestyle{python}{style=base,language=Python}
\lstdefinestyle{xml}{style=base,language=XML,
  morestring=[b]",moredelim=[s][\color{codeKey}]{<}{>}}
\lstdefinestyle{yaml}{style=base,
  keywords={true,false,null},
  comment=[l]{\#},
  moredelim=[l][\color{codeKey}]{:},
}
\lstset{style=base}

% =====================================================================
%  3. THE FLOWCHART VISUAL LANGUAGE
% =====================================================================
%  Every flowchart in this report uses these shapes and only these.
%  The reader learns the alphabet once, in Fig. 1 (the legend), and can
%  then read every later diagram without a caption.
%
%  SHAPES (ISO 5807 / classic flowchart conventions)
%    terminal   rounded    start of a node, end of a node
%    process    rectangle  a computation
%    decision   diamond    a test, two labelled exits
%    io         parallelogram  a ROS topic read or written
%    subroutine rectangle with double side bars: a call into lib/
%    failure    rectangle, red fill: a recovery / degraded branch
%
%  PHASES (the coloured bands behind the shapes)
%    INITIALISATION  blue   everything executed once, before the loop
%    MAIN LOOP       green  the periodic callback / control tick
%    TERMINATION     red    shutdown, resource release, destruction
% ---------------------------------------------------------------------

\tikzset{
  % ---- node shapes ---------------------------------------------------
  fcterm/.style   = {rounded rectangle, draw=nodeTermLine, thick,
                     fill=nodeTerm, minimum width=26mm, minimum height=7mm,
                     align=center, font=\sffamily\scriptsize\bfseries,
                     inner sep=2pt},
  fcproc/.style   = {rectangle, draw=nodeProcLine, thick, fill=nodeProc,
                     minimum width=34mm, minimum height=8mm, align=center,
                     font=\sffamily\scriptsize, inner sep=3pt},
  fcdec/.style    = {diamond, draw=nodeDecLine, thick, fill=nodeDec,
                     aspect=2.2, align=center, inner sep=1pt,
                     font=\sffamily\scriptsize, minimum width=30mm},
  fcio/.style     = {trapezium, trapezium left angle=70,
                     trapezium right angle=110, draw=nodeIOLine, thick,
                     fill=nodeIO, minimum width=30mm, minimum height=7mm,
                     align=center, font=\sffamily\scriptsize\itshape,
                     inner sep=2pt},
  fcsub/.style    = {rectangle, draw=nodeSubLine, thick, fill=nodeSub,
                     minimum width=34mm, minimum height=8mm, align=center,
                     font=\sffamily\scriptsize, inner sep=3pt,
                     path picture={
                       \draw[nodeSubLine,thick]
                         ($(path picture bounding box.north west)+(1.6mm,0)$)
                         -- ($(path picture bounding box.south west)+(1.6mm,0)$);
                       \draw[nodeSubLine,thick]
                         ($(path picture bounding box.north east)-(1.6mm,0)$)
                         -- ($(path picture bounding box.south east)-(1.6mm,0)$);
                     }},
  fcfail/.style   = {rectangle, draw=nodeFailLine, thick, fill=nodeFail,
                     minimum width=34mm, minimum height=8mm, align=center,
                     font=\sffamily\scriptsize, inner sep=3pt},
  fcnote/.style   = {rectangle, draw=gray!60, dashed, fill=nodeNote,
                     align=left, font=\sffamily\tiny\itshape, inner sep=3pt,
                     text width=38mm},
  % ---- edges ---------------------------------------------------------
  fcline/.style   = {-{Stealth[length=2mm]}, thick, draw=nodeProcLine},
  fclab/.style    = {font=\sffamily\tiny\bfseries, inner sep=1.5pt,
                     fill=white, text=nodeProcLine},
  fcback/.style   = {fcline, draw=phaseMainLine, dashed},
  % ---- phase bands (drawn on the background layer) --------------------
  bandinit/.style = {rectangle, rounded corners=2mm, fill=phaseInitFill,
                     draw=phaseInitLine, thick, dash pattern=on 3pt off 2pt},
  bandmain/.style = {rectangle, rounded corners=2mm, fill=phaseMainFill,
                     draw=phaseMainLine, thick, dash pattern=on 3pt off 2pt},
  bandend/.style  = {rectangle, rounded corners=2mm, fill=phaseEndFill,
                     draw=phaseEndLine, thick, dash pattern=on 3pt off 2pt},
  bandlab/.style  = {font=\sffamily\tiny\bfseries, align=left},
}

% \phaseband{style}{label}{colour}{fitted nodes}{name}
%   Draws a coloured band behind the listed nodes and writes the phase
%   name vertically down its left edge. Always used inside a
%   \begin{scope}[on background layer] ... \end{scope}.
\newcommand{\phaseband}[5]{%
  \node[#1, fit=#4, inner sep=3.2mm] (#5) {};
  \node[bandlab, text=#3, rotate=90, anchor=south]
       at ($(#5.west)+(2.0mm,0)$) {#2};
}

% Standard vertical flowchart geometry, so all diagrams line up.
\tikzset{
  fcchain/.style = {node distance=6.5mm and 10mm},
}

% =====================================================================
%  4. SWOT TABLES
% =====================================================================
%  Every algorithm and design decision in this report gets the same
%  four-quadrant table. Internal factors on top (Strengths / Weaknesses),
%  external on the bottom (Opportunities / Threats), which is the
%  standard reading order.
%
%  \swot{title}{strengths}{weaknesses}{opportunities}{threats}
% ---------------------------------------------------------------------
\definecolor{swotS}{HTML}{E4F2E4}
\definecolor{swotW}{HTML}{FBE6E4}
\definecolor{swotO}{HTML}{E3ECF9}
\definecolor{swotT}{HTML}{FDF3D8}

\newcommand{\swotcell}[3]{%
  \cellcolor{#1}\begin{minipage}[t]{\linewidth}%
  \vspace{2pt}\textbf{\footnotesize #2}\par\vspace{2pt}
  \scriptsize #3\vspace{3pt}\end{minipage}}

\newenvironment{swottable}[1]
  {\begin{table}[!t]\caption{SWOT --- #1}\centering\footnotesize
   \renewcommand{\arraystretch}{1.15}
   \begin{tabular}{|p{0.455\columnwidth}|p{0.455\columnwidth}|}\hline}
  {\hline\end{tabular}\end{table}}

% Wide variant for the two-column-spanning SWOT tables.
\newenvironment{swottablewide}[1]
  {\begin{table*}[!t]\caption{SWOT --- #1}\centering\footnotesize
   \renewcommand{\arraystretch}{1.15}
   \begin{tabular}{|p{0.47\textwidth}|p{0.47\textwidth}|}\hline}
  {\hline\end{tabular}\end{table*}}

\usepackage{colortbl}

% =====================================================================
%  5. SMALL SEMANTIC MACROS
% =====================================================================
\newcommand{\topicname}[1]{\texttt{\small /#1}}          % a ROS topic
\newcommand{\nodename}[1]{\texttt{\small #1}}            % a ROS node
\newcommand{\filename}[1]{\texttt{\small #1}}            % a source file
\newcommand{\param}[1]{\texttt{\small #1}}               % a parameter
% \degree is not a LaTeX primitive -- it comes from gensymb or
% siunitx, neither of which is loaded here. Without this the first
% \SI{360}{\degree} kills the build with "Undefined control
% sequence", which is exactly what it did.
\providecommand{\degree}{\ensuremath{^\circ}}
\newcommand{\SI}[2]{#1\,\text{#2}}                       % 0.50 m/s etc.

% Cross-phase reference: "as in Phase A, Sec. X"
\newcommand{\phaseA}{Phase~A}
\newcommand{\phaseB}{Phase~B}
\newcommand{\phaseC}{Phase~C}

% Measured-value emphasis, so every number traceable to a log stands out
\newcommand{\measured}[1]{\textbf{#1}}

% =====================================================================
%  6. MEASURED vs MODELLED
% =====================================================================
%  Phase C reports an energy budget that is DERIVED from a stated model
%  rather than read off an instrument, because the simulator has no
%  battery and the robot has no current sensor. Putting a modelled
%  number in the same typeface as a measured one is how a reader is
%  misled without a single false sentence being written, so the two are
%  marked differently and the distinction is stated once, here, and
%  again where the model is defined (Sec. \ref{sec:phaseC_route}).
%
%    \measured{41.3 m}   read from a log
%    \modelled{5.6 Wh}   computed from assumptions named in the text
% ---------------------------------------------------------------------
\newcommand{\modelled}[1]{\textit{#1}\,$^{\dagger}$}
\newcommand{\modelnote}{\footnotesize$^{\dagger}$~derived from the
  power model of Sec.~\ref{sec:phaseC_route}, whose constants are
  assumptions and not measurements.}

% A station label -- P2,5R -- appears constantly in Phase C.
\newcommand{\station}[1]{\texttt{\small #1}}
